% Options for packages loaded elsewhere
\PassOptionsToPackage{unicode}{hyperref}
\PassOptionsToPackage{hyphens}{url}
\PassOptionsToPackage{dvipsnames,svgnames,x11names}{xcolor}
\documentclass[
  10pt,
  fontset=fandol]{ctexart}
\usepackage{xcolor}
\usepackage[margin=16mm]{geometry}
\usepackage{amsmath,amssymb}
\setcounter{secnumdepth}{-\maxdimen} % remove section numbering
\usepackage{iftex}
\ifPDFTeX
  \usepackage[T1]{fontenc}
  \usepackage[utf8]{inputenc}
  \usepackage{textcomp} % provide euro and other symbols
\else % if luatex or xetex
  \usepackage{unicode-math} % this also loads fontspec
  \defaultfontfeatures{Scale=MatchLowercase}
  \defaultfontfeatures[\rmfamily]{Ligatures=TeX,Scale=1}
\fi
\usepackage{lmodern}
\ifPDFTeX\else
  % xetex/luatex font selection
\fi
% Use upquote if available, for straight quotes in verbatim environments
\IfFileExists{upquote.sty}{\usepackage{upquote}}{}
\IfFileExists{microtype.sty}{% use microtype if available
  \usepackage[]{microtype}
  \UseMicrotypeSet[protrusion]{basicmath} % disable protrusion for tt fonts
}{}
\makeatletter
\@ifundefined{KOMAClassName}{% if non-KOMA class
  \IfFileExists{parskip.sty}{%
    \usepackage{parskip}
  }{% else
    \setlength{\parindent}{0pt}
    \setlength{\parskip}{6pt plus 2pt minus 1pt}}
}{% if KOMA class
  \KOMAoptions{parskip=half}}
\makeatother
\setlength{\emergencystretch}{3em} % prevent overfull lines
\providecommand{\tightlist}{%
  \setlength{\itemsep}{0pt}\setlength{\parskip}{0pt}}
\usepackage{amsmath,amssymb,mathtools,needspace}
\xeCJKDeclareCharClass{CJK}{"2032,"0400->"04FF,"2160->"216B,"2460->"2473,"25A0,"25C7,"25CB,"25CF}
\IfFontExistsTF{Arial Unicode MS}{\xeCJKsetup{AutoFallBack=true}\setCJKfallbackfamilyfont{\CJKrmdefault}{Arial Unicode MS}}{}
\setcounter{secnumdepth}{0}
\setlength{\emergencystretch}{2em}
\allowdisplaybreaks
\usepackage{bookmark}
\IfFileExists{xurl.sty}{\usepackage{xurl}}{} % add URL line breaks if available
\urlstyle{same}
\hypersetup{
  pdftitle={多元最小二乘拟合},
  colorlinks=true,
  linkcolor={Maroon},
  filecolor={Maroon},
  citecolor={Blue},
  urlcolor={Blue},
  pdfcreator={LaTeX via pandoc}}

\title{多元最小二乘拟合}
\author{}
\date{}

\begin{document}
\maketitle

\subsubsection{3.6.3 多元最小二乘拟合}\label{ux591aux5143ux6700ux5c0fux4e8cux4e58ux62dfux5408}

上面介绍的最小二乘法的有关概念与方法可推广到多元函数，例如，已知多元函数

\[
y=f(x_1,x_2,\cdots,x_l)
\]

的一组测量数据 \((x_{1i},x_{2i},\cdots,x_{li},y_i)\)（\(i=1,2,\cdots,m\)），以及一组权系数 \(\omega_i>0\)（\(i=1,2,\cdots,m\)），要求函数

\[
S_n(x_1,x_2,\cdots,x_l)=\sum_{k=1}^na_k\varphi_k(x_1,x_2,\cdots,x_l),\quad n\leqslant m,
\]

使得

\[
F(a_0,a_1,\cdots,a_n)=\sum_{i=1}^m\omega_i[y_i-S_n(x_{1i},x_{2i},\cdots,x_{li})]^2
\]

最小，这与式 (3.6.4) 的极值问题完全一样，系数 \(a_0,a_2,\cdots,a_n\) 同样满足法方程组 (3.6.6)，只是这里

\[
(\varphi_k,\varphi_j)=\sum_{i=1}^m\omega_i\varphi_k(x_{1i},x_{2i},\cdots,x_{li})\varphi_i(x_{1i},x_{2i},\cdots,x_{li}).
\]

求解法方程组 (3.6.6) 就可得到 \(a_k\)（\(k=0,1,\cdots,n\)），从而得到 \(S_n(x_1,x_2,\cdots,x_l)\)。称 \(S_n(x_1,x_2,\cdots,x_l)\) 为函数 \(f(x_1,x_2,\cdots,x_l)\) 的最小二乘拟合。

\begin{center}\rule{0.5\linewidth}{0.5pt}\end{center}

\subsection{相关原页校注（agent 补充）}\label{ux76f8ux5173ux539fux9875ux6821ux6ce8agent-ux8865ux5145}

以下校注来自本节覆盖的原页，可能也涉及同页相邻小节；原文照录与校正说明分开保留。

\subsubsection{原书 PDF 页 84 的校注}\label{ux539fux4e66-pdf-ux9875-84-ux7684ux6821ux6ce8}

\href{../pages/pdf-084.md}{完整原页转录} · \href{../../source-images/pdf-084.jpeg}{原页图像}

校注记录：ch03b-E012, ch03b-E013。

\begin{quote}
校注（agent 补充）：3.6.3 原书的 \(S_n\) 求和从 \(k=1\) 开始，但后续使用 \(F(a_0,a_1,\ldots,a_n)\) 和 \(a_k\;(k=0,\ldots,n)\)；段中系数列表印为 \(a_0,a_2,\ldots,a_n\)；内积公式右端第二因子的原印下标为 \(i\)，而左端为 \(j\)，相容写法应使用 \(\varphi_j\)。这些原书记号不一致均已保留，见 \href{../assets/ch03b/p084-multivariate-formulas.png}{原文裁切}。
\end{quote}

\subsection{本节覆盖原页的校注记录}\label{ux672cux8282ux8986ux76d6ux539fux9875ux7684ux6821ux6ce8ux8bb0ux5f55}

下列记录按原页关联，含同页相邻小节的校注；计算时同时读取上文校注和完整原页。

\begin{itemize}
\tightlist
\item
  \texttt{ch03b-E012}：原书 PDF 页 84
\item
  \texttt{ch03b-E013}：原书 PDF 页 84
\end{itemize}

\end{document}
